\section{Introduction}
\label{sec:introduction}

Thermal models of laser powder bed fusion (LPBF) connect processing conditions to the extent of melting and to the thermal conditions experienced during solidification. Melt-pool geometry identifies the fused region, whereas solidification analysis also uses temperature gradients and material cooling histories. Heat-transfer calculations have been used to relate process inputs to these thermal quantities and to assess grain-structure trends in selected laser and electron-beam powder-bed cases~\cite{Plotkowski2017Rapid}. Spatially and temporally resolved thermography provides complementary observations of temperature profiles and cooling during laser processing~\cite{hooper2018}. Geometry and temperature history thus serve related but distinct purposes: a geometric fit constrains a melting boundary, while the thermal information associated with that boundary depends on the model used to calculate it.

Moving-source conduction models provide this information without resolving the full motion of the liquid or the absorbing surface. Distributed-source formulations remove the near-source singularity of concentrated-source solutions, and nonlinear enthalpy formulations accommodate temperature-dependent properties and latent heat~\cite{vanelsen2007}. Their practical value is illustrated by the IN625 benchmark conducted on the Additive Manufacturing Metrology Testbed (AMMT)~\cite{lane2020}. Using the measured beam profile, Kollmannsberger et al. fitted a scalar absorbed-power parameter that reproduced length more closely than width and depth; a calibrated directional-conductivity model then improved the geometric agreement across the tested conditions~\cite{kollmannsberger2019}. This improvement depended on the material description as well as the source. Their conductivity data ended at 871~$^\circ$C, below the melting range, and they explicitly treated its high-temperature extrapolation as a physical model requiring calibration~\cite{kollmannsberger2019}. The resulting geometry therefore reflects both the incident heating and the assumed transport beyond the available material data.

The thermal information in such a calibrated field requires a separate assessment. The AMMT modeling study restricted the use of its improved model for predicting temperatures inside the melt pool~\cite{kollmannsberger2019}. A later study of bare-plate 316L tracks demonstrated why temperature observations add information to a geometric fit: several Fresnel and accommodation coefficient pairs in a fluid model gave comparable agreement with cross-sectional area and width-to-depth ratio, but predicted different surface-temperature profiles. Two-color measurements distinguished those thermal predictions~\cite{myers2023}. This evidence establishes the importance of the particular observable used to assess a model. In the AMMT measurements, length is obtained from a radiance profile, while width and depth are obtained from metallographic cross sections~\cite{lane2020}. The dimensions constrain different features of the calculated field, and their agreement does not select the material functions that generated its temperature distribution.

The IN625 material functions used in the present study make this dependence concrete. Conductivity determines diffusive heat transport, while sensible specific heat enters the enthalpy stored during heating and released during cooling. The supplier tables provide conductivity only to 982~$^\circ$C and sensible specific heat only to 1093~$^\circ$C, both below the listed 1290--1350~$^\circ$C melting range. The bulletin describes conductivity as measured and specific heat as calculated typical alloy data~\cite{specialmetals625}. Both functions must therefore be extended before the model can evaluate transport and sensible storage throughout the adopted melting interval and above it. Assessments of nickel-superalloy properties also show that liquid-property estimates depend on the measurements and estimation procedure available~\cite{mills2006}. The fitted source-power factor is consequently conditional on the chosen extrapolation functions; a geometric calibration does not supply the missing high-temperature material data.

The extrapolation of sensible properties must be distinguished from the phase-change prescription in the same enthalpy balance. Latent heat is stored and released according to the adopted liquid-fraction relation as well as the total heat of fusion~\cite{vanelsen2007}. A linear fraction across a specified freezing interval is an established continuum approximation~\cite{dynamic2024}. Its endpoints, however, need not coincide across material descriptions: an IN625 model for rapid additive-manufacturing cooling used a lower eutectic endpoint from Gulliver--Scheil calculations, whereas the supplier bulletin lists a melting range~\cite{dynamic2024,specialmetals625}. Changing that interval redistributes latent storage along the temperature axis. Comparing conductivity and sensible-heat-capacity extrapolations with a common phase law therefore gives their responses a defined meaning within the same thermal model.

With the phase thresholds fixed, the rear temperature profile can resolve features that total pool length combines. A solidus-crossing length depends on the front and rear boundary positions. At the rear, the solidus and liquidus crossings describe both the position of the thermal tail relative to the source and the spatial separation across the adopted freezing interval. Displacing the two rear crossings together changes total length while preserving approximately the same separation. Changing their separation alters the spatial interval over which material cools between the adopted thresholds. Conductivity and sensible storage affect the same coupled field, so an overall length response alone cannot distinguish these changes. Resolving the rear pair provides a direct way to interpret which part of the calculated thermal structure responds to the prescribed property functions.

The material time associated with this spatial interval also depends on scan speed. For a steadily translating field, material fixed in the laboratory passes through the rear solidus--liquidus separation in a time equal to that separation divided by the speed. Similar spatial intervals can thus correspond to different passage times and interval-mean cooling rates. Constant-speed conversion from position to time is used in the AMMT cooling analysis, although that analysis considers a below-solidus interval~\cite{lane2020}. The trajectory represented by the conversion remains consequential: pulsed LPBF thermography resolves histories affected by successive heating events~\cite{hooper2018}, and fluid calculations show that liquid motion changes the thermal excursions followed by material parcels~\cite{Hou2024Dissolution}. Passage through a quasi-steady conduction field describes a stationary-material trajectory whose time interpretation is tied to that field and its adopted freezing interval.

Taken together, the need to extrapolate the material inputs and the different information carried by geometry and material time motivate the present assessment of a length-calibrated IN625 conduction field. In a nonlinear moving-frame enthalpy model, an effective source-power fraction is fitted to condition B thermographic length and retained for two complementary comparisons across three AMMT bare-plate conditions. Width, depth and the other operating conditions assess the geometry associated with that fit. At B, independent conductivity and sensible-heat-capacity contrasts compare linear extrapolation using the final tabulated slope with constant endpoint extrapolation under a common phase law. Resolving the rear solidus--liquidus positions and separation distinguishes thermal-tail displacement from a change in interval width, while the operating-condition comparison relates each spatial interval to material passage time at its scan speed. The objective is to determine how the prescribed property functions alter the spatial and temporal interpretation of the calibrated field.
